\documentclass[14pt,a4paper]{extarticle}
\usepackage[T1]{fontenc}
\usepackage[utf8]{inputenc}
\usepackage{amsmath,amssymb}
\usepackage{newtxtext,newtxmath}
\usepackage[a4paper,top=10mm,bottom=15mm,left=13mm,right=13mm,footskip=9mm]{geometry}
\usepackage{xcolor,graphicx,array,tabularx,booktabs,enumitem,fancyhdr}
\usepackage[hidelinks]{hyperref}
\graphicspath{{figures/branding/}}
\definecolor{Brand}{RGB}{0,90,170}
\definecolor{WorkBlue}{HTML}{174CA0}
\definecolor{RuleGray}{HTML}{B9BEC5}
\setlength{\parindent}{0pt}
\setlength{\parskip}{0pt}
\setlength{\abovedisplayskip}{4pt}\setlength{\belowdisplayskip}{4pt}
\setlength{\abovedisplayshortskip}{2pt}\setlength{\belowdisplayshortskip}{2pt}
\raggedbottom
\emergencystretch=1.5em
% ---- editable labels -------------------------------------------------------
\newcommand{\QuizLabel}{Quiz 3}
\newcommand{\TermLabel}{Fall 2026}
\newcommand{\DateLabel}{Thursday, 5 November 2026}
\newcommand{\TimeLabel}{20 minutes}
% ----------------------------------------------------------------------------
\newcommand{\CurVersion}{A}
\newcommand{\CurPage}{1}
\pagestyle{fancy}\fancyhf{}
\renewcommand{\headrulewidth}{0pt}\renewcommand{\footrulewidth}{0.3pt}
\fancyfoot[L]{\fontsize{8}{9.5}\selectfont\begin{tabular}[b]{@{}l@{}}Dr. Yousef Dabboorasad\quad|\quad Math 231 -- \QuizLabel{} (Sections 4.1--4.4)\\ \textcolor{gray}{Based on Anton \& Rorres, \textit{Elementary Linear Algebra}, 12e}\end{tabular}}
\fancyfoot[R]{\fontsize{8.5}{10}\selectfont Version \CurVersion\quad\CurPage/2}
\newcommand{\Logo}[2]{\IfFileExists{figures/branding/#1}{\includegraphics[height=#2]{#1}}{}}
\newcommand{\LogoW}[2]{\IfFileExists{figures/branding/#1}{\includegraphics[width=#2]{#1}}{}}
\newcommand{\HeaderOne}[1]{%
\renewcommand{\CurVersion}{#1}\renewcommand{\CurPage}{1}%
\noindent\begin{minipage}[c]{34mm}\LogoW{math_department_logo.png}{33mm}\end{minipage}%
\hfill\begin{minipage}[c]{\dimexpr\linewidth-56mm\relax}\centering
{\fontsize{10.5}{12}\selectfont\bfseries Qatar University}\\[-1pt]
{\fontsize{9.5}{11}\selectfont Department of Mathematics, Statistics and Physics}\\[2pt]
{\fontsize{16}{18}\selectfont\bfseries Math 231 -- Linear Algebra\quad\textcolor{Brand}{\QuizLabel}}\\[2pt]
{\fontsize{9}{10.5}\selectfont Sections 4.1--4.4 (Real Vector Spaces, Subspaces, Spanning Sets, Linear Independence)}\\[1pt]
{\fontsize{10}{11.5}\selectfont \TermLabel\quad\textbullet\quad \DateLabel\quad\textbullet\quad Time: \TimeLabel\quad\textbullet\quad Total: 10 marks}
\end{minipage}\hfill
\begin{minipage}[c]{18mm}\raggedleft\Logo{qatar_university_logo.png}{14mm}\end{minipage}\par
\vspace{1.5mm}\textcolor{Brand}{\rule{\linewidth}{0.8pt}}\par\vspace{1.5mm}
{\fontsize{11}{13}\selectfont Name: \rule{58mm}{0.3pt}\quad ID: \rule{30mm}{0.3pt}\quad Section: \rule{16mm}{0.3pt}}\hfill
{\fontsize{13}{15}\selectfont\bfseries\textcolor{Brand}{Version #1}}\par\vspace{1mm}
{\fontsize{9}{10.5}\selectfont\color{black!70}Tick exactly one box for each True/False statement and exactly one option for each multiple-choice question; an item with more than one answer receives zero. Show all your work in Problem 3. Programmable calculators and mobile phones are not allowed.}\par\vspace{2mm}
}
\newcommand{\HeaderTwo}[1]{%
\renewcommand{\CurVersion}{#1}\renewcommand{\CurPage}{2}%
\noindent{\fontsize{10.5}{12}\selectfont\bfseries Math 231 -- Linear Algebra\quad\textcolor{Brand}{\QuizLabel}\quad Sections 4.1--4.4}\hfill
{\fontsize{12}{14}\selectfont\bfseries\textcolor{Brand}{Version #1}}\quad{\fontsize{10}{11.5}\selectfont Page 2 of 2}\par
\vspace{1mm}\textcolor{Brand}{\rule{\linewidth}{0.8pt}}\par\vspace{2.5mm}
}
\newcommand{\PTitle}[2]{{\bfseries #1}\hfill{\bfseries [#2]}\par}
\newcommand{\PNote}[1]{{\fontsize{9.5}{11}\selectfont\color{black!70}#1}\par}
\newcommand{\Opt}[2]{$#1$~#2}
\newcommand{\Sol}[1]{{\fontsize{9.5}{11}\selectfont\color{WorkBlue}#1}\par}
\newcommand{\Conclude}[1]{\par\vspace{1pt}\noindent\fcolorbox{WorkBlue}{blue!4}{\textbf{\textcolor{WorkBlue}{#1}}}\par\vspace{2pt}}
\newcommand{\workline}{\noindent\textcolor{RuleGray}{\makebox[\linewidth]{\dotfill}}\par}

\hypersetup{pdftitle={MATH 231 Quiz 3 Sections 4.1-4.4 -- Student paper -- Versions A-J},pdfauthor={Dr. Yousef Dabboorasad},pdfsubject={MATH 231 Quiz 3; Sections 4.1--4.4; versions A--J}}
\begin{document}
\pdfbookmark[0]{Version A}{version-A}
\HeaderOne{A}
\PTitle{Problem 1. True or False}{2 marks}
\PNote{Each statement is worth 0.5 mark. Tick \textbf{True} if the statement is always true, otherwise tick \textbf{False}. No justification is needed.}
\begin{tabularx}{\linewidth}{@{}X@{\hspace{3mm}}>{\centering\arraybackslash}p{21mm}@{}}
 & {\fontsize{9}{10}\selectfont True\hfill False}\\[0.5mm]
\textbf{(a)}~In any vector space $V$, $0\mathbf{u}=\mathbf{0}$ for every vector $\mathbf{u}$ in $V$. & $\square$\hfill$\square$\\[4.2mm]
\textbf{(b)}~The set of all invertible $2\times 2$ matrices is a subspace of $M_{22}$. & $\square$\hfill$\square$\\[4.2mm]
\textbf{(c)}~If $S=\{\mathbf{w}_1,\dots,\mathbf{w}_r\}$ is a set of vectors in a vector space $V$, then $\operatorname{span}(S)$ is a subspace of $V$. & $\square$\hfill$\square$\\[4.2mm]
\textbf{(d)}~If $\{\mathbf{v}_1,\mathbf{v}_2,\mathbf{v}_3\}$ is linearly dependent, then $\mathbf{v}_3$ must be a linear combination of $\mathbf{v}_1$ and $\mathbf{v}_2$. & $\square$\hfill$\square$
\end{tabularx}\par
\vspace{3mm}\textcolor{RuleGray}{\rule{\linewidth}{0.4pt}}\par\vspace{2mm}
\PTitle{Problem 2. Multiple choice}{4 marks}
\PNote{Each question is worth 1 mark. Tick exactly one option.}
\vspace{0.8mm}
\textbf{(a)}~Let $V=\mathbb{R}^2$ with the usual addition $(x_1,x_2)+(y_1,y_2)=(x_1+y_1,\;x_2+y_2)$ and the scalar multiplication $k(x_1,x_2)=(kx_1,\;0)$. With these operations $V$ is \emph{not} a vector space: exactly one of the following four axioms fails. Which one?\par\vspace{0.8mm}
\begin{tabularx}{\linewidth}{@{}X@{\hspace{3mm}}X@{}}
$\square$ A.~Axiom 7: $k(\mathbf{u}+\mathbf{v})=k\mathbf{u}+k\mathbf{v}$ & $\square$ B.~Axiom 10: $1\mathbf{u}=\mathbf{u}$\\[0.8mm]
$\square$ C.~Axiom 8: $(k+m)\mathbf{u}=k\mathbf{u}+m\mathbf{u}$ & $\square$ D.~Axiom 9: $k(m\mathbf{u})=(km)\mathbf{u}$
\end{tabularx}\par
\vspace{2.6mm}
\textbf{(b)}~Which of the following subsets $W$ is a \textbf{subspace} of $V=\mathbb{R}^2$ (with the usual operations)?\par\vspace{0.8mm}
\begin{tabularx}{\linewidth}{@{}X@{\hspace{3mm}}X@{}}
$\square$ A.~$W=\{(x,y) : x+y=1\}$ & $\square$ B.~$W=\{(x,y) : y=x^2\}$\\[0.8mm]
$\square$ C.~$W=\{(x,y) : y=3x\}$ & $\square$ D.~$W=\{(x,y) : x\ge0,\ y\ge0\}$
\end{tabularx}\par
\vspace{2.6mm}
\textbf{(c)}~Let $\mathbf{v}_1=(1,0,2)$ and $\mathbf{v}_2=(0,1,-1)$ in $\mathbb{R}^3$. Which of the following vectors belongs to $\operatorname{span}\{\mathbf{v}_1,\mathbf{v}_2\}$?\par\vspace{0.8mm}
\begin{tabular*}{\linewidth}{@{}l@{\extracolsep{\fill}}lll@{}}
$\square$ A.~$(3,1,7)$ & $\square$ B.~$(3,1,5)$ & $\square$ C.~$(1,1,2)$ & $\square$ D.~$(2,2,0)$\end{tabular*}\par
\vspace{2.6mm}
\textbf{(d)}~Let $f_1=x$ and $f_2=e^{x}$ be functions in $C^{\infty}(-\infty,\infty)$. Compute the Wronskian $W(x)$ of $f_1,f_2$ and choose the correct statement.\par\vspace{0.8mm}
\begin{tabularx}{\linewidth}{@{}X@{\hspace{3mm}}X@{}}
$\square$ A.~$W(x)=(x-1)e^{x}$; the set is linearly dependent because $W(1)=0$ & $\square$ B.~$W(x)=(x-1)e^{x}$; the set is linearly independent\\[0.8mm]
$\square$ C.~$W(x)=(x+1)e^{x}$; the set is linearly independent & $\square$ D.~$W(x)=0$ for all $x$; the Wronskian test is inconclusive
\end{tabularx}\par
\newpage

\pdfbookmark[1]{Version A: Problem 3}{written-A}
\HeaderTwo{A}
\PTitle{Problem 3. Linear independence and span}{4 marks}
\vspace{1.5mm}Let $S=\{\mathbf{v}_1,\mathbf{v}_2,\mathbf{v}_3\}$ be the set of vectors in $\mathbb{R}^3$ given by\[\mathbf{v}_1=(1,2,-1),\qquad \mathbf{v}_2=(2,1,1),\qquad \mathbf{v}_3=(5,4,1),\]and let $\mathbf{u}=(1,1,3)$.
\begin{enumerate}[label=\textbf{(\alph*)},leftmargin=9mm,itemsep=1.5mm,topsep=2mm]
\item \textbf{[2.5 marks]} Determine whether $S$ is linearly independent or linearly dependent. Show all your work.
\item \textbf{[1.5 marks]} Determine whether $\mathbf{u}$ belongs to $\operatorname{span}(S)$, that is, whether $\mathbf{u}$ is a linear combination of $\mathbf{v}_1,\mathbf{v}_2,\mathbf{v}_3$. If it is, write $\mathbf{u}$ as such a linear combination; if it is not, explain why. (You may reuse the row operations from part (a).)
\end{enumerate}
\vspace{2mm}\textbf{Working:}\par\vspace{5.5mm}
\workline\vspace{6.6mm}
\workline\vspace{6.6mm}
\workline\vspace{6.6mm}
\workline\vspace{6.6mm}
\workline\vspace{6.6mm}
\workline\vspace{6.6mm}
\workline\vspace{6.6mm}
\workline\vspace{6.6mm}
\workline\vspace{6.6mm}
\workline\vspace{6.6mm}
\workline\vspace{6.6mm}
\workline
\newpage
\pdfbookmark[0]{Version B}{version-B}
\HeaderOne{B}
\PTitle{Problem 1. True or False}{2 marks}
\PNote{Each statement is worth 0.5 mark. Tick \textbf{True} if the statement is always true, otherwise tick \textbf{False}. No justification is needed.}
\begin{tabularx}{\linewidth}{@{}X@{\hspace{3mm}}>{\centering\arraybackslash}p{21mm}@{}}
 & {\fontsize{9}{10}\selectfont True\hfill False}\\[0.5mm]
\textbf{(a)}~A set with two operations that satisfies nine of the ten vector space axioms is still a vector space. & $\square$\hfill$\square$\\[4.2mm]
\textbf{(b)}~If $A$ is an $m\times n$ matrix, then the solution set of the homogeneous system $A\mathbf{x}=\mathbf{0}$ is a subspace of $\mathbb{R}^n$. & $\square$\hfill$\square$\\[4.2mm]
\textbf{(c)}~$\operatorname{span}\{1,\,x,\,x^2\}=P_3$. & $\square$\hfill$\square$\\[4.2mm]
\textbf{(d)}~Every set of four vectors in $\mathbb{R}^3$ is linearly dependent. & $\square$\hfill$\square$
\end{tabularx}\par
\vspace{3mm}\textcolor{RuleGray}{\rule{\linewidth}{0.4pt}}\par\vspace{2mm}
\PTitle{Problem 2. Multiple choice}{4 marks}
\PNote{Each question is worth 1 mark. Tick exactly one option.}
\vspace{0.8mm}
\textbf{(a)}~Let $V=\mathbb{R}^2$ with the usual addition $(x_1,x_2)+(y_1,y_2)=(x_1+y_1,\;x_2+y_2)$ and the scalar multiplication $k(x_1,x_2)=(kx_1,\;x_2)$. With these operations $V$ is \emph{not} a vector space: exactly one of the following four axioms fails. Which one?\par\vspace{0.8mm}
\begin{tabularx}{\linewidth}{@{}X@{\hspace{3mm}}X@{}}
$\square$ A.~Axiom 10: $1\mathbf{u}=\mathbf{u}$ & $\square$ B.~Axiom 7: $k(\mathbf{u}+\mathbf{v})=k\mathbf{u}+k\mathbf{v}$\\[0.8mm]
$\square$ C.~Axiom 8: $(k+m)\mathbf{u}=k\mathbf{u}+m\mathbf{u}$ & $\square$ D.~Axiom 6: $k\mathbf{u}$ is in $V$ (closure under scalar multiplication)
\end{tabularx}\par
\vspace{2.6mm}
\textbf{(b)}~Which of the following subsets of $V=M_{22}$ (with the usual operations) is \textbf{not} a subspace of $V$?\par\vspace{0.8mm}
\begin{tabularx}{\linewidth}{@{}X@{\hspace{3mm}}X@{}}
$\square$ A.~all symmetric matrices ($A^T=A$) & $\square$ B.~all invertible matrices\\[0.8mm]
$\square$ C.~all matrices with $\operatorname{tr}(A)=0$ & $\square$ D.~all matrices $\begin{bmatrix}a&b\\ c&d\end{bmatrix}$ with $a=0$
\end{tabularx}\par
\vspace{2.6mm}
\textbf{(c)}~Let $p_1=1+x^{2}$ and $p_2=x$ in $P_2$. Which of the following polynomials belongs to $\operatorname{span}\{p_1,p_2\}$?\par\vspace{0.8mm}
\begin{tabular*}{\linewidth}{@{}l@{\extracolsep{\fill}}lll@{}}
$\square$ A.~$3+2x-3x^{2}$ & $\square$ B.~$2+3x+3x^{2}$ & $\square$ C.~$3+2x+3x^{2}$ & $\square$ D.~$3+3x+2x^{2}$\end{tabular*}\par
\vspace{2.6mm}
\textbf{(d)}~Let $f_1=e^{x}$ and $f_2=e^{2x}$ be functions in $C^{\infty}(-\infty,\infty)$. Compute the Wronskian $W(x)$ of $f_1,f_2$ and choose the correct statement.\par\vspace{0.8mm}
\begin{tabularx}{\linewidth}{@{}X@{\hspace{3mm}}X@{}}
$\square$ A.~$W(x)=3e^{3x}$; the set is linearly independent & $\square$ B.~$W(x)=0$ for all $x$; the Wronskian test is inconclusive\\[0.8mm]
$\square$ C.~$W(x)=e^{3x}$; the set is linearly dependent because $e^{2x}$ is a multiple of $e^{x}$ & $\square$ D.~$W(x)=e^{3x}$; the set is linearly independent
\end{tabularx}\par
\newpage

\pdfbookmark[1]{Version B: Problem 3}{written-B}
\HeaderTwo{B}
\PTitle{Problem 3. Linear independence and span}{4 marks}
\vspace{1.5mm}Let $S=\{p_1,p_2,p_3\}$ be the set of polynomials in $P_2$ given by\[p_1=1+x+x^{2},\qquad p_2=1-x+2x^{2},\qquad p_3=2+x-x^{2},\]and let $u=3-2x+2x^{2}$.
\begin{enumerate}[label=\textbf{(\alph*)},leftmargin=9mm,itemsep=1.5mm,topsep=2mm]
\item \textbf{[2.5 marks]} Determine whether $S$ is linearly independent or linearly dependent. Show all your work.
\item \textbf{[1.5 marks]} Determine whether $u$ belongs to $\operatorname{span}(S)$, that is, whether $u$ is a linear combination of $p_1,p_2,p_3$. If it is, write $u$ as such a linear combination; if it is not, explain why. (You may reuse the row operations from part (a).)
\end{enumerate}
\vspace{2mm}\textbf{Working:}\par\vspace{5.5mm}
\workline\vspace{6.6mm}
\workline\vspace{6.6mm}
\workline\vspace{6.6mm}
\workline\vspace{6.6mm}
\workline\vspace{6.6mm}
\workline\vspace{6.6mm}
\workline\vspace{6.6mm}
\workline\vspace{6.6mm}
\workline\vspace{6.6mm}
\workline\vspace{6.6mm}
\workline\vspace{6.6mm}
\workline
\newpage
\pdfbookmark[0]{Version C}{version-C}
\HeaderOne{C}
\PTitle{Problem 1. True or False}{2 marks}
\PNote{Each statement is worth 0.5 mark. Tick \textbf{True} if the statement is always true, otherwise tick \textbf{False}. No justification is needed.}
\begin{tabularx}{\linewidth}{@{}X@{\hspace{3mm}}>{\centering\arraybackslash}p{21mm}@{}}
 & {\fontsize{9}{10}\selectfont True\hfill False}\\[0.5mm]
\textbf{(a)}~If $\mathbf{u}$ is a vector in a vector space and $k$ is a scalar with $k\mathbf{u}=\mathbf{0}$, then $k=0$ or $\mathbf{u}=\mathbf{0}$. & $\square$\hfill$\square$\\[4.2mm]
\textbf{(b)}~If $W_1$ and $W_2$ are subspaces of a vector space $V$, then $W_1\cup W_2$ is always a subspace of $V$. & $\square$\hfill$\square$\\[4.2mm]
\textbf{(c)}~The vectors $(1,0,0)$ and $(0,1,0)$ span $\mathbb{R}^3$. & $\square$\hfill$\square$\\[4.2mm]
\textbf{(d)}~A set with exactly two vectors is linearly independent if and only if neither vector is a scalar multiple of the other. & $\square$\hfill$\square$
\end{tabularx}\par
\vspace{3mm}\textcolor{RuleGray}{\rule{\linewidth}{0.4pt}}\par\vspace{2mm}
\PTitle{Problem 2. Multiple choice}{4 marks}
\PNote{Each question is worth 1 mark. Tick exactly one option.}
\vspace{0.8mm}
\textbf{(a)}~Let $V=\mathbb{R}^2$ with the usual addition $(x_1,x_2)+(y_1,y_2)=(x_1+y_1,\;x_2+y_2)$ and the scalar multiplication $k(x_1,x_2)=(k^2x_1,\;k^2x_2)$. With these operations $V$ is \emph{not} a vector space: exactly one of the following four axioms fails. Which one?\par\vspace{0.8mm}
\begin{tabularx}{\linewidth}{@{}X@{\hspace{3mm}}X@{}}
$\square$ A.~Axiom 8: $(k+m)\mathbf{u}=k\mathbf{u}+m\mathbf{u}$ & $\square$ B.~Axiom 9: $k(m\mathbf{u})=(km)\mathbf{u}$\\[0.8mm]
$\square$ C.~Axiom 10: $1\mathbf{u}=\mathbf{u}$ & $\square$ D.~Axiom 7: $k(\mathbf{u}+\mathbf{v})=k\mathbf{u}+k\mathbf{v}$
\end{tabularx}\par
\vspace{2.6mm}
\textbf{(b)}~Which of the following subsets $W$ is a \textbf{subspace} of $V=P_2$ (with the usual operations)?\par\vspace{0.8mm}
\begin{tabularx}{\linewidth}{@{}X@{\hspace{3mm}}X@{}}
$\square$ A.~$W=\{p\in P_2 : \deg(p)=2\}$ & $\square$ B.~$W=\{a+bx+cx^2 : a,b,c \text{ are integers}\}$\\[0.8mm]
$\square$ C.~$W=\{p\in P_2 : p(0)=1\}$ & $\square$ D.~$W=\{a+bx+cx^2 : a+b+c=0\}$
\end{tabularx}\par
\vspace{2.6mm}
\textbf{(c)}~Let $\mathbf{v}_1=(1,0,-1)$ and $\mathbf{v}_2=(0,1,3)$ in $\mathbb{R}^3$. Which of the following vectors belongs to $\operatorname{span}\{\mathbf{v}_1,\mathbf{v}_2\}$?\par\vspace{0.8mm}
\begin{tabular*}{\linewidth}{@{}l@{\extracolsep{\fill}}lll@{}}
$\square$ A.~$(2,1,1)$ & $\square$ B.~$(2,1,5)$ & $\square$ C.~$(1,2,-5)$ & $\square$ D.~$(2,1,-1)$\end{tabular*}\par
\vspace{2.6mm}
\textbf{(d)}~Let $f_1=1$ and $f_2=e^{-x}$ be functions in $C^{\infty}(-\infty,\infty)$. Compute the Wronskian $W(x)$ of $f_1,f_2$ and choose the correct statement.\par\vspace{0.8mm}
\begin{tabularx}{\linewidth}{@{}X@{\hspace{3mm}}X@{}}
$\square$ A.~$W(x)=-e^{-x}$; the set is linearly independent & $\square$ B.~$W(x)=e^{-x}$; the set is linearly independent\\[0.8mm]
$\square$ C.~$W(x)=-e^{-x}$; the set is linearly dependent because $W(x)<0$ & $\square$ D.~$W(x)=0$ for all $x$; the Wronskian test is inconclusive
\end{tabularx}\par
\newpage

\pdfbookmark[1]{Version C: Problem 3}{written-C}
\HeaderTwo{C}
\PTitle{Problem 3. Linear independence and span}{4 marks}
\vspace{1.5mm}Let $S=\{\mathbf{v}_1,\mathbf{v}_2,\mathbf{v}_3\}$ be the set of vectors in $\mathbb{R}^3$ given by\[\mathbf{v}_1=(1,1,2),\qquad \mathbf{v}_2=(1,0,1),\qquad \mathbf{v}_3=(2,1,1),\]and let $\mathbf{u}=(3,3,4)$.
\begin{enumerate}[label=\textbf{(\alph*)},leftmargin=9mm,itemsep=1.5mm,topsep=2mm]
\item \textbf{[2.5 marks]} Determine whether $S$ is linearly independent or linearly dependent. Show all your work.
\item \textbf{[1.5 marks]} Determine whether $\mathbf{u}$ belongs to $\operatorname{span}(S)$, that is, whether $\mathbf{u}$ is a linear combination of $\mathbf{v}_1,\mathbf{v}_2,\mathbf{v}_3$. If it is, write $\mathbf{u}$ as such a linear combination; if it is not, explain why. (You may reuse the row operations from part (a).)
\end{enumerate}
\vspace{2mm}\textbf{Working:}\par\vspace{5.5mm}
\workline\vspace{6.6mm}
\workline\vspace{6.6mm}
\workline\vspace{6.6mm}
\workline\vspace{6.6mm}
\workline\vspace{6.6mm}
\workline\vspace{6.6mm}
\workline\vspace{6.6mm}
\workline\vspace{6.6mm}
\workline\vspace{6.6mm}
\workline\vspace{6.6mm}
\workline\vspace{6.6mm}
\workline
\newpage
\pdfbookmark[0]{Version D}{version-D}
\HeaderOne{D}
\PTitle{Problem 1. True or False}{2 marks}
\PNote{Each statement is worth 0.5 mark. Tick \textbf{True} if the statement is always true, otherwise tick \textbf{False}. No justification is needed.}
\begin{tabularx}{\linewidth}{@{}X@{\hspace{3mm}}>{\centering\arraybackslash}p{21mm}@{}}
 & {\fontsize{9}{10}\selectfont True\hfill False}\\[0.5mm]
\textbf{(a)}~In a vector space, if $k\mathbf{u}=\mathbf{0}$ for a scalar $k$ and a vector $\mathbf{u}$, then both $k=0$ and $\mathbf{u}=\mathbf{0}$. & $\square$\hfill$\square$\\[4.2mm]
\textbf{(b)}~If $W_1$ and $W_2$ are subspaces of a vector space $V$, then $W_1\cap W_2$ is a subspace of $V$. & $\square$\hfill$\square$\\[4.2mm]
\textbf{(c)}~Every vector in a set $S$ belongs to $\operatorname{span}(S)$. & $\square$\hfill$\square$\\[4.2mm]
\textbf{(d)}~The functions $e^{x}$ and $6e^{x}$ form a linearly independent set in $C(-\infty,\infty)$. & $\square$\hfill$\square$
\end{tabularx}\par
\vspace{3mm}\textcolor{RuleGray}{\rule{\linewidth}{0.4pt}}\par\vspace{2mm}
\PTitle{Problem 2. Multiple choice}{4 marks}
\PNote{Each question is worth 1 mark. Tick exactly one option.}
\vspace{0.8mm}
\textbf{(a)}~Let $V=\mathbb{R}^2$ with the usual addition $(x_1,x_2)+(y_1,y_2)=(x_1+y_1,\;x_2+y_2)$ and the scalar multiplication $k(x_1,x_2)=(0,\;kx_2)$. With these operations $V$ is \emph{not} a vector space: exactly one of the following four axioms fails. Which one?\par\vspace{0.8mm}
\begin{tabularx}{\linewidth}{@{}X@{\hspace{3mm}}X@{}}
$\square$ A.~Axiom 6: $k\mathbf{u}$ is in $V$ (closure under scalar multiplication) & $\square$ B.~Axiom 4: there is a zero vector $\mathbf{0}$ with $\mathbf{0}+\mathbf{u}=\mathbf{u}$\\[0.8mm]
$\square$ C.~Axiom 7: $k(\mathbf{u}+\mathbf{v})=k\mathbf{u}+k\mathbf{v}$ & $\square$ D.~Axiom 10: $1\mathbf{u}=\mathbf{u}$
\end{tabularx}\par
\vspace{2.6mm}
\textbf{(b)}~Which of the following subsets of $V=\mathbb{R}^3$ (with the usual operations) is \textbf{not} a subspace of $V$?\par\vspace{0.8mm}
\begin{tabularx}{\linewidth}{@{}X@{\hspace{3mm}}X@{}}
$\square$ A.~$W=\{(x,y,z) : x+y+z=1\}$ & $\square$ B.~$W=\{(x,y,z) : x+y+z=0\}$\\[0.8mm]
$\square$ C.~$W=\{(x,y,z) : x=y=z\}$ & $\square$ D.~$W=\{(x,y,z) : x=2y \text{ and } z=0\}$
\end{tabularx}\par
\vspace{2.6mm}
\textbf{(c)}~Let $p_1=1+x$ and $p_2=x^{2}$ in $P_2$. Which of the following polynomials belongs to $\operatorname{span}\{p_1,p_2\}$?\par\vspace{0.8mm}
\begin{tabular*}{\linewidth}{@{}l@{\extracolsep{\fill}}lll@{}}
$\square$ A.~$2-2x-x^{2}$ & $\square$ B.~$2+x-x^{2}$ & $\square$ C.~$1+2x+2x^{2}$ & $\square$ D.~$2+2x-x^{2}$\end{tabular*}\par
\vspace{2.6mm}
\textbf{(d)}~Let $f_1=x$ and $f_2=x^{2}$ be functions in $C^{\infty}(-\infty,\infty)$. Compute the Wronskian $W(x)$ of $f_1,f_2$ and choose the correct statement.\par\vspace{0.8mm}
\begin{tabularx}{\linewidth}{@{}X@{\hspace{3mm}}X@{}}
$\square$ A.~$W(x)=x^{2}$; the set is linearly dependent because $W(0)=0$ & $\square$ B.~$W(x)=x^{2}$; the Wronskian test is inconclusive because $W(0)=0$\\[0.8mm]
$\square$ C.~$W(x)=x^{2}$; the set is linearly independent & $\square$ D.~$W(x)=3x^{2}$; the set is linearly independent
\end{tabularx}\par
\newpage

\pdfbookmark[1]{Version D: Problem 3}{written-D}
\HeaderTwo{D}
\PTitle{Problem 3. Linear independence and span}{4 marks}
\vspace{1.5mm}Let $S=\{p_1,p_2,p_3\}$ be the set of polynomials in $P_2$ given by\[p_1=1+x,\qquad p_2=1-x+x^{2},\qquad p_3=1+3x-x^{2},\]and let $u=1+2x+x^{2}$.
\begin{enumerate}[label=\textbf{(\alph*)},leftmargin=9mm,itemsep=1.5mm,topsep=2mm]
\item \textbf{[2.5 marks]} Determine whether $S$ is linearly independent or linearly dependent. Show all your work.
\item \textbf{[1.5 marks]} Determine whether $u$ belongs to $\operatorname{span}(S)$, that is, whether $u$ is a linear combination of $p_1,p_2,p_3$. If it is, write $u$ as such a linear combination; if it is not, explain why. (You may reuse the row operations from part (a).)
\end{enumerate}
\vspace{2mm}\textbf{Working:}\par\vspace{5.5mm}
\workline\vspace{6.6mm}
\workline\vspace{6.6mm}
\workline\vspace{6.6mm}
\workline\vspace{6.6mm}
\workline\vspace{6.6mm}
\workline\vspace{6.6mm}
\workline\vspace{6.6mm}
\workline\vspace{6.6mm}
\workline\vspace{6.6mm}
\workline\vspace{6.6mm}
\workline\vspace{6.6mm}
\workline
\newpage
\pdfbookmark[0]{Version E}{version-E}
\HeaderOne{E}
\PTitle{Problem 1. True or False}{2 marks}
\PNote{Each statement is worth 0.5 mark. Tick \textbf{True} if the statement is always true, otherwise tick \textbf{False}. No justification is needed.}
\begin{tabularx}{\linewidth}{@{}X@{\hspace{3mm}}>{\centering\arraybackslash}p{21mm}@{}}
 & {\fontsize{9}{10}\selectfont True\hfill False}\\[0.5mm]
\textbf{(a)}~The set $F(-\infty,\infty)$ of all real-valued functions defined on $(-\infty,\infty)$, with pointwise addition and scalar multiplication, is a vector space. & $\square$\hfill$\square$\\[4.2mm]
\textbf{(b)}~The set of all $n\times n$ matrices $A$ with $\operatorname{tr}(A)=0$ is a subspace of $M_{nn}$. & $\square$\hfill$\square$\\[4.2mm]
\textbf{(c)}~If $\operatorname{span}(S)=\mathbb{R}^3$, then $S$ must contain exactly three vectors. & $\square$\hfill$\square$\\[4.2mm]
\textbf{(d)}~If the Wronskian $W(x)$ of two functions satisfies $W(0)=0$, then the Wronskian test is inconclusive. & $\square$\hfill$\square$
\end{tabularx}\par
\vspace{3mm}\textcolor{RuleGray}{\rule{\linewidth}{0.4pt}}\par\vspace{2mm}
\PTitle{Problem 2. Multiple choice}{4 marks}
\PNote{Each question is worth 1 mark. Tick exactly one option.}
\vspace{0.8mm}
\textbf{(a)}~Let $V=\mathbb{R}^3$ with the usual addition $(x_1,x_2,x_3)+(y_1,y_2,y_3)=(x_1+y_1,\;x_2+y_2,\;x_3+y_3)$ and the scalar multiplication $k(x_1,x_2,x_3)=(kx_1,\;x_2,\;x_3)$. With these operations $V$ is \emph{not} a vector space: exactly one of the following four axioms fails. Which one?\par\vspace{0.8mm}
\begin{tabularx}{\linewidth}{@{}X@{\hspace{3mm}}X@{}}
$\square$ A.~Axiom 9: $k(m\mathbf{u})=(km)\mathbf{u}$ & $\square$ B.~Axiom 8: $(k+m)\mathbf{u}=k\mathbf{u}+m\mathbf{u}$\\[0.8mm]
$\square$ C.~Axiom 10: $1\mathbf{u}=\mathbf{u}$ & $\square$ D.~Axiom 2: $\mathbf{u}+\mathbf{v}=\mathbf{v}+\mathbf{u}$
\end{tabularx}\par
\vspace{2.6mm}
\textbf{(b)}~Which of the following subsets $W$ is a \textbf{subspace} of $V=F(-\infty,\infty)$ (with the usual operations)?\par\vspace{0.8mm}
\begin{tabularx}{\linewidth}{@{}X@{\hspace{3mm}}X@{}}
$\square$ A.~$W=\{f : f(0)=1\}$ & $\square$ B.~$W=\{f : f(x)\ge0 \text{ for all } x\}$\\[0.8mm]
$\square$ C.~$W=C(-\infty,\infty)$, all continuous functions & $\square$ D.~$W=\{f : f(1)=2\}$
\end{tabularx}\par
\vspace{2.6mm}
\textbf{(c)}~Let $\mathbf{v}_1=(1,1,0)$ and $\mathbf{v}_2=(0,1,1)$ in $\mathbb{R}^3$. Which of the following vectors belongs to $\operatorname{span}\{\mathbf{v}_1,\mathbf{v}_2\}$?\par\vspace{0.8mm}
\begin{tabular*}{\linewidth}{@{}l@{\extracolsep{\fill}}lll@{}}
$\square$ A.~$(2,3,5)$ & $\square$ B.~$(2,5,3)$ & $\square$ C.~$(2,5,4)$ & $\square$ D.~$(1,1,1)$\end{tabular*}\par
\vspace{2.6mm}
\textbf{(d)}~Let $f_1=e^{x}$ and $f_2=xe^{x}$ be functions in $C^{\infty}(-\infty,\infty)$. Compute the Wronskian $W(x)$ of $f_1,f_2$ and choose the correct statement.\par\vspace{0.8mm}
\begin{tabularx}{\linewidth}{@{}X@{\hspace{3mm}}X@{}}
$\square$ A.~$W(x)=e^{2x}$; the set is linearly independent & $\square$ B.~$W(x)=xe^{2x}$; the set is linearly independent\\[0.8mm]
$\square$ C.~$W(x)=xe^{2x}$; the set is linearly dependent because $W(0)=0$ & $\square$ D.~$W(x)=e^{2x}$; the set is linearly dependent because $xe^{x}$ is a multiple of $e^{x}$
\end{tabularx}\par
\newpage

\pdfbookmark[1]{Version E: Problem 3}{written-E}
\HeaderTwo{E}
\PTitle{Problem 3. Linear independence and span}{4 marks}
\vspace{1.5mm}Let $S=\{\mathbf{v}_1,\mathbf{v}_2,\mathbf{v}_3\}$ be the set of vectors in $\mathbb{R}^3$ given by\[\mathbf{v}_1=(1,-1,2),\qquad \mathbf{v}_2=(1,1,0),\qquad \mathbf{v}_3=(1,-3,4),\]and let $\mathbf{u}=(1,2,1)$.
\begin{enumerate}[label=\textbf{(\alph*)},leftmargin=9mm,itemsep=1.5mm,topsep=2mm]
\item \textbf{[2.5 marks]} Determine whether $S$ is linearly independent or linearly dependent. Show all your work.
\item \textbf{[1.5 marks]} Determine whether $\mathbf{u}$ belongs to $\operatorname{span}(S)$, that is, whether $\mathbf{u}$ is a linear combination of $\mathbf{v}_1,\mathbf{v}_2,\mathbf{v}_3$. If it is, write $\mathbf{u}$ as such a linear combination; if it is not, explain why. (You may reuse the row operations from part (a).)
\end{enumerate}
\vspace{2mm}\textbf{Working:}\par\vspace{5.5mm}
\workline\vspace{6.6mm}
\workline\vspace{6.6mm}
\workline\vspace{6.6mm}
\workline\vspace{6.6mm}
\workline\vspace{6.6mm}
\workline\vspace{6.6mm}
\workline\vspace{6.6mm}
\workline\vspace{6.6mm}
\workline\vspace{6.6mm}
\workline\vspace{6.6mm}
\workline\vspace{6.6mm}
\workline
\newpage
\pdfbookmark[0]{Version F}{version-F}
\HeaderOne{F}
\PTitle{Problem 1. True or False}{2 marks}
\PNote{Each statement is worth 0.5 mark. Tick \textbf{True} if the statement is always true, otherwise tick \textbf{False}. No justification is needed.}
\begin{tabularx}{\linewidth}{@{}X@{\hspace{3mm}}>{\centering\arraybackslash}p{21mm}@{}}
 & {\fontsize{9}{10}\selectfont True\hfill False}\\[0.5mm]
\textbf{(a)}~To show that a set with two operations is a vector space, it is enough to verify Axiom 1 (closure under addition) and Axiom 6 (closure under scalar multiplication). & $\square$\hfill$\square$\\[4.2mm]
\textbf{(b)}~The first quadrant $W=\{(x,y)\in\mathbb{R}^2 : x\ge 0,\ y\ge 0\}$ is a subspace of $\mathbb{R}^2$. & $\square$\hfill$\square$\\[4.2mm]
\textbf{(c)}~$\operatorname{span}\{1,\,x,\,x^2\}=P_2$. & $\square$\hfill$\square$\\[4.2mm]
\textbf{(d)}~If the Wronskian of $f_1,\dots,f_n$ is not identically zero on $(-\infty,\infty)$, then the functions are linearly independent. & $\square$\hfill$\square$
\end{tabularx}\par
\vspace{3mm}\textcolor{RuleGray}{\rule{\linewidth}{0.4pt}}\par\vspace{2mm}
\PTitle{Problem 2. Multiple choice}{4 marks}
\PNote{Each question is worth 1 mark. Tick exactly one option.}
\vspace{0.8mm}
\textbf{(a)}~Let $V=\mathbb{R}^3$ with the usual addition $(x_1,x_2,x_3)+(y_1,y_2,y_3)=(x_1+y_1,\;x_2+y_2,\;x_3+y_3)$ and the scalar multiplication $k(x_1,x_2,x_3)=(kx_1,\;kx_2,\;0)$. With these operations $V$ is \emph{not} a vector space: exactly one of the following four axioms fails. Which one?\par\vspace{0.8mm}
\begin{tabularx}{\linewidth}{@{}X@{\hspace{3mm}}X@{}}
$\square$ A.~Axiom 1: $\mathbf{u}+\mathbf{v}$ is in $V$ (closure under addition) & $\square$ B.~Axiom 8: $(k+m)\mathbf{u}=k\mathbf{u}+m\mathbf{u}$\\[0.8mm]
$\square$ C.~Axiom 9: $k(m\mathbf{u})=(km)\mathbf{u}$ & $\square$ D.~Axiom 10: $1\mathbf{u}=\mathbf{u}$
\end{tabularx}\par
\vspace{2.6mm}
\textbf{(b)}~Which of the following subsets $W$ is a \textbf{subspace} of $V=M_{22}$ (with the usual operations)?\par\vspace{0.8mm}
\begin{tabularx}{\linewidth}{@{}X@{\hspace{3mm}}X@{}}
$\square$ A.~all matrices with $\det(A)=0$ & $\square$ B.~all matrices with $\operatorname{tr}(A)=0$\\[0.8mm]
$\square$ C.~all matrices of the form $\begin{bmatrix}1&a\\ b&c\end{bmatrix}$ & $\square$ D.~all matrices whose four entries are all $\ge0$
\end{tabularx}\par
\vspace{2.6mm}
\textbf{(c)}~Let $p_1=1-x$ and $p_2=1+x^{2}$ in $P_2$. Which of the following polynomials belongs to $\operatorname{span}\{p_1,p_2\}$?\par\vspace{0.8mm}
\begin{tabular*}{\linewidth}{@{}l@{\extracolsep{\fill}}lll@{}}
$\square$ A.~$1+2x-3x^{2}$ & $\square$ B.~$3+2x+x^{2}$ & $\square$ C.~$1+2x+3x^{2}$ & $\square$ D.~$1-2x+3x^{2}$\end{tabular*}\par
\vspace{2.6mm}
\textbf{(d)}~Let $f_1=1$ and $f_2=x^{2}$ be functions in $C^{\infty}(-\infty,\infty)$. Compute the Wronskian $W(x)$ of $f_1,f_2$ and choose the correct statement.\par\vspace{0.8mm}
\begin{tabularx}{\linewidth}{@{}X@{\hspace{3mm}}X@{}}
$\square$ A.~$W(x)=2x$; the set is linearly dependent because $W(0)=0$ & $\square$ B.~$W(x)=-2x$; the set is linearly independent\\[0.8mm]
$\square$ C.~$W(x)=2x$; the Wronskian test is inconclusive because $W(0)=0$ & $\square$ D.~$W(x)=2x$; the set is linearly independent
\end{tabularx}\par
\newpage

\pdfbookmark[1]{Version F: Problem 3}{written-F}
\HeaderTwo{F}
\PTitle{Problem 3. Linear independence and span}{4 marks}
\vspace{1.5mm}Let $S=\{p_1,p_2,p_3\}$ be the set of polynomials in $P_2$ given by\[p_1=1+x^{2},\qquad p_2=x+x^{2},\qquad p_3=1-x+x^{2},\]and let $u=-1+4x+x^{2}$.
\begin{enumerate}[label=\textbf{(\alph*)},leftmargin=9mm,itemsep=1.5mm,topsep=2mm]
\item \textbf{[2.5 marks]} Determine whether $S$ is linearly independent or linearly dependent. Show all your work.
\item \textbf{[1.5 marks]} Determine whether $u$ belongs to $\operatorname{span}(S)$, that is, whether $u$ is a linear combination of $p_1,p_2,p_3$. If it is, write $u$ as such a linear combination; if it is not, explain why. (You may reuse the row operations from part (a).)
\end{enumerate}
\vspace{2mm}\textbf{Working:}\par\vspace{5.5mm}
\workline\vspace{6.6mm}
\workline\vspace{6.6mm}
\workline\vspace{6.6mm}
\workline\vspace{6.6mm}
\workline\vspace{6.6mm}
\workline\vspace{6.6mm}
\workline\vspace{6.6mm}
\workline\vspace{6.6mm}
\workline\vspace{6.6mm}
\workline\vspace{6.6mm}
\workline\vspace{6.6mm}
\workline
\newpage
\pdfbookmark[0]{Version G}{version-G}
\HeaderOne{G}
\PTitle{Problem 1. True or False}{2 marks}
\PNote{Each statement is worth 0.5 mark. Tick \textbf{True} if the statement is always true, otherwise tick \textbf{False}. No justification is needed.}
\begin{tabularx}{\linewidth}{@{}X@{\hspace{3mm}}>{\centering\arraybackslash}p{21mm}@{}}
 & {\fontsize{9}{10}\selectfont True\hfill False}\\[0.5mm]
\textbf{(a)}~In a vector space $V$, $(-1)\mathbf{u}=-\mathbf{u}$ for every vector $\mathbf{u}$ in $V$. & $\square$\hfill$\square$\\[4.2mm]
\textbf{(b)}~The set $W=\{a+bx+cx^2 : a,\,b,\,c \text{ are integers}\}$ is a subspace of $P_2$. & $\square$\hfill$\square$\\[4.2mm]
\textbf{(c)}~The zero vector belongs to $\operatorname{span}(S)$ for every nonempty set $S$ of vectors in $V$. & $\square$\hfill$\square$\\[4.2mm]
\textbf{(d)}~If the Wronskian of $f_1,\dots,f_n$ is identically zero on $(-\infty,\infty)$, then the functions are linearly dependent. & $\square$\hfill$\square$
\end{tabularx}\par
\vspace{3mm}\textcolor{RuleGray}{\rule{\linewidth}{0.4pt}}\par\vspace{2mm}
\PTitle{Problem 2. Multiple choice}{4 marks}
\PNote{Each question is worth 1 mark. Tick exactly one option.}
\vspace{0.8mm}
\textbf{(a)}~Let $V=\mathbb{R}^2$ with the usual addition $(x_1,x_2)+(y_1,y_2)=(x_1+y_1,\;x_2+y_2)$ and the scalar multiplication $k(x_1,x_2)=(kx_1,\;k^2x_2)$. With these operations $V$ is \emph{not} a vector space: exactly one of the following four axioms fails. Which one?\par\vspace{0.8mm}
\begin{tabularx}{\linewidth}{@{}X@{\hspace{3mm}}X@{}}
$\square$ A.~Axiom 10: $1\mathbf{u}=\mathbf{u}$ & $\square$ B.~Axiom 9: $k(m\mathbf{u})=(km)\mathbf{u}$\\[0.8mm]
$\square$ C.~Axiom 8: $(k+m)\mathbf{u}=k\mathbf{u}+m\mathbf{u}$ & $\square$ D.~Axiom 7: $k(\mathbf{u}+\mathbf{v})=k\mathbf{u}+k\mathbf{v}$
\end{tabularx}\par
\vspace{2.6mm}
\textbf{(b)}~Which of the following subsets of $V=\mathbb{R}^2$ (with the usual operations) is \textbf{not} a subspace of $V$?\par\vspace{0.8mm}
\begin{tabularx}{\linewidth}{@{}X@{\hspace{3mm}}X@{}}
$\square$ A.~$W=\{(x,y) : x=0\}$ (the $y$-axis) & $\square$ B.~$W=\{(x,y) : 2x-5y=0\}$\\[0.8mm]
$\square$ C.~$W=\{(0,0)\}$ & $\square$ D.~$W=\{(x,y) : xy\ge0\}$
\end{tabularx}\par
\vspace{2.6mm}
\textbf{(c)}~Let $\mathbf{v}_1=(1,0,0)$ and $\mathbf{v}_2=(0,1,2)$ in $\mathbb{R}^3$. Which of the following vectors belongs to $\operatorname{span}\{\mathbf{v}_1,\mathbf{v}_2\}$?\par\vspace{0.8mm}
\begin{tabular*}{\linewidth}{@{}l@{\extracolsep{\fill}}lll@{}}
$\square$ A.~$(4,2,1)$ & $\square$ B.~$(1,2,3)$ & $\square$ C.~$(0,1,3)$ & $\square$ D.~$(4,1,2)$\end{tabular*}\par
\vspace{2.6mm}
\textbf{(d)}~Let $f_1=e^{2x}$ and $f_2=e^{3x}$ be functions in $C^{\infty}(-\infty,\infty)$. Compute the Wronskian $W(x)$ of $f_1,f_2$ and choose the correct statement.\par\vspace{0.8mm}
\begin{tabularx}{\linewidth}{@{}X@{\hspace{3mm}}X@{}}
$\square$ A.~$W(x)=5e^{5x}$; the set is linearly independent & $\square$ B.~$W(x)=e^{5x}$; the set is linearly independent\\[0.8mm]
$\square$ C.~$W(x)=e^{5x}$; the set is linearly dependent because $e^{3x}=e^{x}\cdot e^{2x}$ & $\square$ D.~$W(x)=e^{6x}$; the set is linearly independent
\end{tabularx}\par
\newpage

\pdfbookmark[1]{Version G: Problem 3}{written-G}
\HeaderTwo{G}
\PTitle{Problem 3. Linear independence and span}{4 marks}
\vspace{1.5mm}Let $S=\{\mathbf{v}_1,\mathbf{v}_2,\mathbf{v}_3\}$ be the set of vectors in $\mathbb{R}^3$ given by\[\mathbf{v}_1=(1,2,1),\qquad \mathbf{v}_2=(2,1,-1),\qquad \mathbf{v}_3=(1,1,2),\]and let $\mathbf{u}=(4,3,-3)$.
\begin{enumerate}[label=\textbf{(\alph*)},leftmargin=9mm,itemsep=1.5mm,topsep=2mm]
\item \textbf{[2.5 marks]} Determine whether $S$ is linearly independent or linearly dependent. Show all your work.
\item \textbf{[1.5 marks]} Determine whether $\mathbf{u}$ belongs to $\operatorname{span}(S)$, that is, whether $\mathbf{u}$ is a linear combination of $\mathbf{v}_1,\mathbf{v}_2,\mathbf{v}_3$. If it is, write $\mathbf{u}$ as such a linear combination; if it is not, explain why. (You may reuse the row operations from part (a).)
\end{enumerate}
\vspace{2mm}\textbf{Working:}\par\vspace{5.5mm}
\workline\vspace{6.6mm}
\workline\vspace{6.6mm}
\workline\vspace{6.6mm}
\workline\vspace{6.6mm}
\workline\vspace{6.6mm}
\workline\vspace{6.6mm}
\workline\vspace{6.6mm}
\workline\vspace{6.6mm}
\workline\vspace{6.6mm}
\workline\vspace{6.6mm}
\workline\vspace{6.6mm}
\workline
\newpage
\pdfbookmark[0]{Version H}{version-H}
\HeaderOne{H}
\PTitle{Problem 1. True or False}{2 marks}
\PNote{Each statement is worth 0.5 mark. Tick \textbf{True} if the statement is always true, otherwise tick \textbf{False}. No justification is needed.}
\begin{tabularx}{\linewidth}{@{}X@{\hspace{3mm}}>{\centering\arraybackslash}p{21mm}@{}}
 & {\fontsize{9}{10}\selectfont True\hfill False}\\[0.5mm]
\textbf{(a)}~If $\mathbf{u}$ is a vector in a vector space and $k$ is a scalar with $k\mathbf{u}=\mathbf{0}$, then necessarily $k=0$. & $\square$\hfill$\square$\\[4.2mm]
\textbf{(b)}~The set $W=\{a+bx+cx^2 : a+b+c=0\}$ is a subspace of $P_2$. & $\square$\hfill$\square$\\[4.2mm]
\textbf{(c)}~The set $S=\{(1,-1,0),\,(1,2,1),\,(1,0,1),\,(3,3,-1)\}$ cannot span $\mathbb{R}^3$ because it contains more than three vectors. & $\square$\hfill$\square$\\[4.2mm]
\textbf{(d)}~Every set of five matrices in $M_{22}$ is linearly dependent. & $\square$\hfill$\square$
\end{tabularx}\par
\vspace{3mm}\textcolor{RuleGray}{\rule{\linewidth}{0.4pt}}\par\vspace{2mm}
\PTitle{Problem 2. Multiple choice}{4 marks}
\PNote{Each question is worth 1 mark. Tick exactly one option.}
\vspace{0.8mm}
\textbf{(a)}~Let $V=\mathbb{R}^3$ with the usual addition $(x_1,x_2,x_3)+(y_1,y_2,y_3)=(x_1+y_1,\;x_2+y_2,\;x_3+y_3)$ and the scalar multiplication $k(x_1,x_2,x_3)=(0,\;kx_2,\;kx_3)$. With these operations $V$ is \emph{not} a vector space: exactly one of the following four axioms fails. Which one?\par\vspace{0.8mm}
\begin{tabularx}{\linewidth}{@{}X@{\hspace{3mm}}X@{}}
$\square$ A.~Axiom 10: $1\mathbf{u}=\mathbf{u}$ & $\square$ B.~Axiom 8: $(k+m)\mathbf{u}=k\mathbf{u}+m\mathbf{u}$\\[0.8mm]
$\square$ C.~Axiom 4: there is a zero vector $\mathbf{0}$ with $\mathbf{0}+\mathbf{u}=\mathbf{u}$ & $\square$ D.~Axiom 1: $\mathbf{u}+\mathbf{v}$ is in $V$ (closure under addition)
\end{tabularx}\par
\vspace{2.6mm}
\textbf{(b)}~Which of the following subsets of $V=F(-\infty,\infty)$ (with the usual operations) is \textbf{not} a subspace of $V$?\par\vspace{0.8mm}
\begin{tabularx}{\linewidth}{@{}X@{\hspace{3mm}}X@{}}
$\square$ A.~$W=\{f : f(0)=1\}$ & $\square$ B.~$W=\{f : f(0)=0\}$\\[0.8mm]
$\square$ C.~$W=C^1(-\infty,\infty)$, all continuously differentiable functions & $\square$ D.~$W=P_2$, all polynomials of degree at most 2
\end{tabularx}\par
\vspace{2.6mm}
\textbf{(c)}~Let $p_1=x+x^{2}$ and $p_2=1$ in $P_2$. Which of the following polynomials belongs to $\operatorname{span}\{p_1,p_2\}$?\par\vspace{0.8mm}
\begin{tabular*}{\linewidth}{@{}l@{\extracolsep{\fill}}lll@{}}
$\square$ A.~$5+2x+2x^{2}$ & $\square$ B.~$5+2x+5x^{2}$ & $\square$ C.~$2+5x+2x^{2}$ & $\square$ D.~$5+2x-2x^{2}$\end{tabular*}\par
\vspace{2.6mm}
\textbf{(d)}~Let $f_1=x$ and $f_2=1+x$ be functions in $C^{\infty}(-\infty,\infty)$. Compute the Wronskian $W(x)$ of $f_1,f_2$ and choose the correct statement.\par\vspace{0.8mm}
\begin{tabularx}{\linewidth}{@{}X@{\hspace{3mm}}X@{}}
$\square$ A.~$W(x)=1$; the set is linearly independent & $\square$ B.~$W(x)=-1$; the set is linearly dependent because $W(x)$ is negative\\[0.8mm]
$\square$ C.~$W(x)=-1$; the set is linearly independent & $\square$ D.~$W(x)=0$ for all $x$; the Wronskian test is inconclusive
\end{tabularx}\par
\newpage

\pdfbookmark[1]{Version H: Problem 3}{written-H}
\HeaderTwo{H}
\PTitle{Problem 3. Linear independence and span}{4 marks}
\vspace{1.5mm}Let $S=\{p_1,p_2,p_3\}$ be the set of polynomials in $P_2$ given by\[p_1=2+x-x^{2},\qquad p_2=1+x+x^{2},\qquad p_3=4+3x+x^{2},\]and let $u=1+2x+x^{2}$.
\begin{enumerate}[label=\textbf{(\alph*)},leftmargin=9mm,itemsep=1.5mm,topsep=2mm]
\item \textbf{[2.5 marks]} Determine whether $S$ is linearly independent or linearly dependent. Show all your work.
\item \textbf{[1.5 marks]} Determine whether $u$ belongs to $\operatorname{span}(S)$, that is, whether $u$ is a linear combination of $p_1,p_2,p_3$. If it is, write $u$ as such a linear combination; if it is not, explain why. (You may reuse the row operations from part (a).)
\end{enumerate}
\vspace{2mm}\textbf{Working:}\par\vspace{5.5mm}
\workline\vspace{6.6mm}
\workline\vspace{6.6mm}
\workline\vspace{6.6mm}
\workline\vspace{6.6mm}
\workline\vspace{6.6mm}
\workline\vspace{6.6mm}
\workline\vspace{6.6mm}
\workline\vspace{6.6mm}
\workline\vspace{6.6mm}
\workline\vspace{6.6mm}
\workline\vspace{6.6mm}
\workline
\newpage
\pdfbookmark[0]{Version I}{version-I}
\HeaderOne{I}
\PTitle{Problem 1. True or False}{2 marks}
\PNote{Each statement is worth 0.5 mark. Tick \textbf{True} if the statement is always true, otherwise tick \textbf{False}. No justification is needed.}
\begin{tabularx}{\linewidth}{@{}X@{\hspace{3mm}}>{\centering\arraybackslash}p{21mm}@{}}
 & {\fontsize{9}{10}\selectfont True\hfill False}\\[0.5mm]
\textbf{(a)}~In a vector space $V$, the equation $k\mathbf{0}=\mathbf{0}$ holds only when the scalar $k$ equals $0$. & $\square$\hfill$\square$\\[4.2mm]
\textbf{(b)}~The set $C(-\infty,\infty)$ of all continuous functions on $(-\infty,\infty)$ is a subspace of $F(-\infty,\infty)$. & $\square$\hfill$\square$\\[4.2mm]
\textbf{(c)}~If $\mathbf{u}\in\operatorname{span}\{\mathbf{v}_1,\mathbf{v}_2\}$, then $\mathbf{u}=k_1\mathbf{v}_1+k_2\mathbf{v}_2$ for some scalars $k_1,k_2$. & $\square$\hfill$\square$\\[4.2mm]
\textbf{(d)}~Every set of two nonzero vectors in $\mathbb{R}^3$ is linearly independent. & $\square$\hfill$\square$
\end{tabularx}\par
\vspace{3mm}\textcolor{RuleGray}{\rule{\linewidth}{0.4pt}}\par\vspace{2mm}
\PTitle{Problem 2. Multiple choice}{4 marks}
\PNote{Each question is worth 1 mark. Tick exactly one option.}
\vspace{0.8mm}
\textbf{(a)}~Let $V=\mathbb{R}^2$ with the usual addition $(x_1,x_2)+(y_1,y_2)=(x_1+y_1,\;x_2+y_2)$ and the scalar multiplication $k(x_1,x_2)=(x_1,\;kx_2)$. With these operations $V$ is \emph{not} a vector space: exactly one of the following four axioms fails. Which one?\par\vspace{0.8mm}
\begin{tabularx}{\linewidth}{@{}X@{\hspace{3mm}}X@{}}
$\square$ A.~Axiom 7: $k(\mathbf{u}+\mathbf{v})=k\mathbf{u}+k\mathbf{v}$ & $\square$ B.~Axiom 9: $k(m\mathbf{u})=(km)\mathbf{u}$\\[0.8mm]
$\square$ C.~Axiom 8: $(k+m)\mathbf{u}=k\mathbf{u}+m\mathbf{u}$ & $\square$ D.~Axiom 10: $1\mathbf{u}=\mathbf{u}$
\end{tabularx}\par
\vspace{2.6mm}
\textbf{(b)}~Which of the following subsets $W$ is a \textbf{subspace} of $V=\mathbb{R}^3$ (with the usual operations)?\par\vspace{0.8mm}
\begin{tabularx}{\linewidth}{@{}X@{\hspace{3mm}}X@{}}
$\square$ A.~$W=\{(x,y,z) : x\ge0\}$ & $\square$ B.~$W=\{(x,y,z) : x^2+y^2+z^2\le1\}$\\[0.8mm]
$\square$ C.~$W=\{(x,y,z) : xyz=0\}$ & $\square$ D.~$W=\{(x,y,z) : z=x-2y\}$
\end{tabularx}\par
\vspace{2.6mm}
\textbf{(c)}~Let $\mathbf{v}_1=(1,0,1)$ and $\mathbf{v}_2=(0,1,1)$ in $\mathbb{R}^3$. Which of the following vectors belongs to $\operatorname{span}\{\mathbf{v}_1,\mathbf{v}_2\}$?\par\vspace{0.8mm}
\begin{tabular*}{\linewidth}{@{}l@{\extracolsep{\fill}}lll@{}}
$\square$ A.~$(3,-1,4)$ & $\square$ B.~$(3,1,2)$ & $\square$ C.~$(3,-1,2)$ & $\square$ D.~$(1,1,1)$\end{tabular*}\par
\vspace{2.6mm}
\textbf{(d)}~Let $f_1=x^{2}$ and $f_2=x^{3}$ be functions in $C^{\infty}(-\infty,\infty)$. Compute the Wronskian $W(x)$ of $f_1,f_2$ and choose the correct statement.\par\vspace{0.8mm}
\begin{tabularx}{\linewidth}{@{}X@{\hspace{3mm}}X@{}}
$\square$ A.~$W(x)=x^{4}$; the set is linearly dependent because $W(0)=0$ & $\square$ B.~$W(x)=x^{4}$; the set is linearly independent\\[0.8mm]
$\square$ C.~$W(x)=5x^{4}$; the set is linearly independent & $\square$ D.~$W(x)=0$ for all $x$; the Wronskian test is inconclusive
\end{tabularx}\par
\newpage

\pdfbookmark[1]{Version I: Problem 3}{written-I}
\HeaderTwo{I}
\PTitle{Problem 3. Linear independence and span}{4 marks}
\vspace{1.5mm}Let $S=\{\mathbf{v}_1,\mathbf{v}_2,\mathbf{v}_3\}$ be the set of vectors in $\mathbb{R}^3$ given by\[\mathbf{v}_1=(1,-2,1),\qquad \mathbf{v}_2=(3,1,2),\qquad \mathbf{v}_3=(5,-3,4),\]and let $\mathbf{u}=(2,1,3)$.
\begin{enumerate}[label=\textbf{(\alph*)},leftmargin=9mm,itemsep=1.5mm,topsep=2mm]
\item \textbf{[2.5 marks]} Determine whether $S$ is linearly independent or linearly dependent. Show all your work.
\item \textbf{[1.5 marks]} Determine whether $\mathbf{u}$ belongs to $\operatorname{span}(S)$, that is, whether $\mathbf{u}$ is a linear combination of $\mathbf{v}_1,\mathbf{v}_2,\mathbf{v}_3$. If it is, write $\mathbf{u}$ as such a linear combination; if it is not, explain why. (You may reuse the row operations from part (a).)
\end{enumerate}
\vspace{2mm}\textbf{Working:}\par\vspace{5.5mm}
\workline\vspace{6.6mm}
\workline\vspace{6.6mm}
\workline\vspace{6.6mm}
\workline\vspace{6.6mm}
\workline\vspace{6.6mm}
\workline\vspace{6.6mm}
\workline\vspace{6.6mm}
\workline\vspace{6.6mm}
\workline\vspace{6.6mm}
\workline\vspace{6.6mm}
\workline\vspace{6.6mm}
\workline
\newpage
\pdfbookmark[0]{Version J}{version-J}
\HeaderOne{J}
\PTitle{Problem 1. True or False}{2 marks}
\PNote{Each statement is worth 0.5 mark. Tick \textbf{True} if the statement is always true, otherwise tick \textbf{False}. No justification is needed.}
\begin{tabularx}{\linewidth}{@{}X@{\hspace{3mm}}>{\centering\arraybackslash}p{21mm}@{}}
 & {\fontsize{9}{10}\selectfont True\hfill False}\\[0.5mm]
\textbf{(a)}~In a real vector space, the scalars may be any real numbers, including irrational numbers such as $\sqrt{2}$. & $\square$\hfill$\square$\\[4.2mm]
\textbf{(b)}~If $A$ is an $m\times n$ matrix and $\mathbf{b}\neq\mathbf{0}$, then the solution set of $A\mathbf{x}=\mathbf{b}$ is a subspace of $\mathbb{R}^n$. & $\square$\hfill$\square$\\[4.2mm]
\textbf{(c)}~The polynomial $x^2$ belongs to $\operatorname{span}\{1+x^2,\;x\}$. & $\square$\hfill$\square$\\[4.2mm]
\textbf{(d)}~A set consisting of a single nonzero vector is linearly independent. & $\square$\hfill$\square$
\end{tabularx}\par
\vspace{3mm}\textcolor{RuleGray}{\rule{\linewidth}{0.4pt}}\par\vspace{2mm}
\PTitle{Problem 2. Multiple choice}{4 marks}
\PNote{Each question is worth 1 mark. Tick exactly one option.}
\vspace{0.8mm}
\textbf{(a)}~Let $V=\mathbb{R}^3$ with the usual addition $(x_1,x_2,x_3)+(y_1,y_2,y_3)=(x_1+y_1,\;x_2+y_2,\;x_3+y_3)$ and the scalar multiplication $k(x_1,x_2,x_3)=(x_1,\;x_2,\;kx_3)$. With these operations $V$ is \emph{not} a vector space: exactly one of the following four axioms fails. Which one?\par\vspace{0.8mm}
\begin{tabularx}{\linewidth}{@{}X@{\hspace{3mm}}X@{}}
$\square$ A.~Axiom 6: $k\mathbf{u}$ is in $V$ (closure under scalar multiplication) & $\square$ B.~Axiom 8: $(k+m)\mathbf{u}=k\mathbf{u}+m\mathbf{u}$\\[0.8mm]
$\square$ C.~Axiom 10: $1\mathbf{u}=\mathbf{u}$ & $\square$ D.~Axiom 9: $k(m\mathbf{u})=(km)\mathbf{u}$
\end{tabularx}\par
\vspace{2.6mm}
\textbf{(b)}~Which of the following subsets of $V=P_2$ (with the usual operations) is \textbf{not} a subspace of $V$?\par\vspace{0.8mm}
\begin{tabularx}{\linewidth}{@{}X@{\hspace{3mm}}X@{}}
$\square$ A.~$W=\{p\in P_2 : p(0)=0\}$ & $\square$ B.~$W=\{a+bx+cx^2 : a+b+c=2\}$\\[0.8mm]
$\square$ C.~$W=\{a+bx+cx^2 : a=2c\}$ & $\square$ D.~$W=\{p\in P_2 : p(1)=0\}$
\end{tabularx}\par
\vspace{2.6mm}
\textbf{(c)}~Let $p_1=1+2x$ and $p_2=x^{2}$ in $P_2$. Which of the following polynomials belongs to $\operatorname{span}\{p_1,p_2\}$?\par\vspace{0.8mm}
\begin{tabular*}{\linewidth}{@{}l@{\extracolsep{\fill}}lll@{}}
$\square$ A.~$2+2x-x^{2}$ & $\square$ B.~$2+4x-x^{2}$ & $\square$ C.~$4+2x-x^{2}$ & $\square$ D.~$1+4x-x^{2}$\end{tabular*}\par
\vspace{2.6mm}
\textbf{(d)}~Let $f_1=e^{-x}$ and $f_2=e^{x}$ be functions in $C^{\infty}(-\infty,\infty)$. Compute the Wronskian $W(x)$ of $f_1,f_2$ and choose the correct statement.\par\vspace{0.8mm}
\begin{tabularx}{\linewidth}{@{}X@{\hspace{3mm}}X@{}}
$\square$ A.~$W(x)=2$; the set is linearly independent & $\square$ B.~$W(x)=0$ for all $x$; the Wronskian test is inconclusive\\[0.8mm]
$\square$ C.~$W(x)=-2$; the set is linearly independent & $\square$ D.~$W(x)=2$; the set is linearly dependent because $W(x)$ is constant
\end{tabularx}\par
\newpage

\pdfbookmark[1]{Version J: Problem 3}{written-J}
\HeaderTwo{J}
\PTitle{Problem 3. Linear independence and span}{4 marks}
\vspace{1.5mm}Let $S=\{p_1,p_2,p_3\}$ be the set of polynomials in $P_2$ given by\[p_1=1+2x,\qquad p_2=x-x^{2},\qquad p_3=1+x^{2},\]and let $u=1-x+x^{2}$.
\begin{enumerate}[label=\textbf{(\alph*)},leftmargin=9mm,itemsep=1.5mm,topsep=2mm]
\item \textbf{[2.5 marks]} Determine whether $S$ is linearly independent or linearly dependent. Show all your work.
\item \textbf{[1.5 marks]} Determine whether $u$ belongs to $\operatorname{span}(S)$, that is, whether $u$ is a linear combination of $p_1,p_2,p_3$. If it is, write $u$ as such a linear combination; if it is not, explain why. (You may reuse the row operations from part (a).)
\end{enumerate}
\vspace{2mm}\textbf{Working:}\par\vspace{5.5mm}
\workline\vspace{6.6mm}
\workline\vspace{6.6mm}
\workline\vspace{6.6mm}
\workline\vspace{6.6mm}
\workline\vspace{6.6mm}
\workline\vspace{6.6mm}
\workline\vspace{6.6mm}
\workline\vspace{6.6mm}
\workline\vspace{6.6mm}
\workline\vspace{6.6mm}
\workline\vspace{6.6mm}
\workline

\end{document}